\documentclass[10pt,a4paper]{amsart}
\usepackage[margin=19mm]{geometry}
\usepackage{amsmath,amssymb,mathtools}
\usepackage[scr=rsfs,cal=euler]{mathalfa}
\usepackage{xfrac}
\usepackage[unicode,hidelinks,bookmarks=false]{hyperref}
\newcommand{\ie}{i.e.\ }
\newcommand{\cf}{cf.\ }
\newcommand{\resp}{respectively\ }
\newcommand{\Subsection}{\subsection}
\providecommand{\nobreakdash}{\nobreak\hbox{-}}
\setlength{\emergencystretch}{2em}
\multlinegap0pt
\usepackage{bm}
\usepackage{bbm}
\usepackage{dsfont}
\usepackage{xspace}
\usepackage{cancel}
\usepackage{mathtools}
\usepackage{amsthm}
\makeatletter
\@ifundefined{theorem}{\newtheorem{theorem}{Theorem}}{}
\@ifundefined{lemma}{\newtheorem{lemma}[theorem]{Lemma}}{}
\makeatother
\title[Markov Properties of $k$-Record Processes via Order Statisti]{Markov Properties of $k$-Record Processes via Order Statistics}
\author{Rodrigo Labouriau}
\date{Source: arXiv:2607.11283v2 (CC BY 4.0)}
\begin{document}
\maketitle

\noindent\emph{Source and licence.} Markov Properties of $k$-Record Processes via Order Statistics. Version arXiv:2607.11283v2 (CC BY 4.0), distributed under the Creative Commons Attribution 4.0 International licence, as recorded in the arXiv licence metadata for this exact version and shown on the abstract page https://arxiv.org/abs/2607.11283v2. Full rights notice: licenses/arxiv-2607.11283-CC-BY-4.0.txt.\par\smallskip

\noindent\textit{Editorial scope.} Counted unit: the Theorem whose source label is \texttt{thm:transition-kernel-Un} in the article (source0.tex lines 436--454, source label \texttt{thm:transition-kernel-Un}), delivered together with the article's own complete original proof of it (source0.tex lines 456--538, one \texttt{proof} environment of 83 lines). No mathematics is written, repaired or re-derived: statement and proof are the article's own text line for line. Required context retained uncounted: lem:conditional-upper-order-stats (lines 382--393). Every definition, display and label of those lines is retained unchanged. Omitted from this package and not redistributed: 1--381, 394--435, 455--455, 539--1108. Nothing inside a retained excerpt is omitted or altered: every retained body is delivered exactly as the article's lines are written, so no omission receipt of any kind accompanies this package. Citations: the retained excerpts carry no citation command at all, so no bibliography entry of the article is transcribed and no reference is redistributed. Reference adaptation: none was needed and none was made; every reference inside the delivered unit resolves inside it, so no unresolved reference or citation is produced. Printed slips kept literal: no printed word, symbol, hypothesis or number of the counted statement or of its proof was altered. Layout and class: the article's own class is \texttt{article} with packages including amsmath, amssymb, amsthm, mathtools, enumitem, hyperref, setspace, fancyhdr; the wrapper is the lane's pinned \texttt{amsart} editorial wrapper at 10pt on A4, which restates the theorem-like environments the retained text needs and adds the article's own macro definitions that the retained text needs (none), with extra wrapper packages (bm, bbm, dsfont, xspace, cancel, mathtools). The delivered numbering is the unit's own. Assets: no retained span contains a figure environment, a graphics-inclusion command, a PDF-page inclusion, a table or a TikZ picture; no image file is copied into this package, the package declares no asset dependency and no image file is redistributed with it. Licence: version arXiv:2607.11283v2 of this article is distributed under the Creative Commons Attribution 4.0 International licence, as recorded in the arXiv licence metadata for this exact version and shown on the abstract page https://arxiv.org/abs/2607.11283v2. A full rights notice is delivered at licenses/arxiv-2607.11283-CC-BY-4.0.txt. Characters: every retained excerpt is pure ASCII, so the delivered engine, XeLaTeX with its default Latin Modern text font, reports no missing character.
\begin{lemma}
\label{lem:conditional-upper-order-stats}
Fix $n\ge k$, and let
\[
F_x(y)=\frac{F(y)-F(x)}{1-F(x)}, \qquad y\ge x,
\]
for every $x$ such that $F(x)<1$. Then, conditionally on the event $\{U_n=x\}$, the random vector
\[
\bigl(X_{n-k+2:n},\dots,X_{n:n}\bigr)
\]
has the same distribution as the order statistics of $k-1$ independent random variables with distribution function $F_x$.
\end{lemma}

\begin{theorem}
\label{thm:transition-kernel-Un}
The process $(U_n)_{n\ge k}$ is a time-homogeneous Markov chain 
% with respect to the natural filtration $(\mathcal{F}_n)_{n\ge1}$. 
with respect to its natural filtration
$\mathcal{G}_n=\sigma(U_k,\dots,U_n)$, $n\ge k$.
Its transition kernel $Q_k$ is given by
\begin{equation}
\label{eq:transition-kernel-Un}
Q_k(x,A)
=
F(x)\,\delta_x(A)
+
\int_{A\cap(x,\infty)}
k\frac{(1-F(y))^{k-1}}{(1-F(x))^{k-1}}\,dF(y),
\qquad A\in\mathcal{B}(\mathbb{R}),
\end{equation}
for every $x$ such that $F(x)<1$, where $\delta_x$ denotes the Dirac measure at $x$.
\end{theorem}

\begin{proof}
Fix $n\ge k$ and let $x$ be such that $F(x)<1$. On the event $\{U_n=x\}$ there are exactly $k-1$ sample points strictly larger than $x$, namely
\[
X_{n-k+2:n},\dots,X_{n:n},
\]
and all remaining sample points are less than or equal to $x$.

Let $X_{n+1}$ be the new observation. Then:
\begin{itemize}
\item if $X_{n+1}\le x$, the running $k$th largest value remains unchanged, so $U_{n+1}=x$;
\item if $X_{n+1}>x$, then $U_{n+1}$ is the minimum of the $k$ values
\[
X_{n-k+2:n},\dots,X_{n:n},X_{n+1}.
\]
\end{itemize}

It follows that, for every $y>x$,
\begin{align*}
\mathbb{P}(U_{n+1}>y\mid U_n=x)
&=
\mathbb{P}\bigl(X_{n+1}>y,\ X_{n-k+2:n}>y \mid U_n=x\bigr).
\end{align*}
By independence of $X_{n+1}$ and $\mathcal{F}_n$,
\[
\mathbb{P}(X_{n+1}>y\mid U_n=x)=1-F(y).
\]
Moreover, by Lemma~\ref{lem:conditional-upper-order-stats}, conditionally on $U_n=x$, the variable $X_{n-k+2:n}$ is the minimum of $k-1$ independent random variables with distribution function $F_x$. Therefore
\[
\mathbb{P}(X_{n-k+2:n}>y\mid U_n=x)
=
\left(1-F_x(y)\right)^{k-1}
=
\left(\frac{1-F(y)}{1-F(x)}\right)^{k-1}.
\]
Hence, for $y>x$,
\begin{equation}
\label{eq:tail-kernel-Un}
\mathbb{P}(U_{n+1}>y\mid U_n=x)
=
(1-F(y))
\left(\frac{1-F(y)}{1-F(x)}\right)^{k-1}
=
\frac{(1-F(y))^k}{(1-F(x))^{k-1}}.
\end{equation}

Taking $y\downarrow x$ in \eqref{eq:tail-kernel-Un}, and using continuity of $F$, we obtain
\[
\mathbb{P}(U_{n+1}>x\mid U_n=x)=1-F(x),
\]
so that
\[
\mathbb{P}(U_{n+1}=x\mid U_n=x)=F(x).
\]
Thus the kernel has an atom of mass $F(x)$ at the point $x$.

Finally, for $y>x$,
\[
\mathbb{P}(x<U_{n+1}\le y\mid U_n=x)
=
\frac{(1-F(x))^k-(1-F(y))^k}{(1-F(x))^{k-1}}
=
\int_{(x,y]}
k\frac{(1-F(t))^{k-1}}{(1-F(x))^{k-1}}\,dF(t).
\]
% Together with the atom at $x$, this proves \eqref{eq:transition-kernel-Un}. 
% Since the conditional law of $U_{n+1}$ given $\mathcal{F}_n$ depends on the past only through $U_n$, 
% the process $(U_n)_{n\ge k}$ is a time-homogeneous Markov chain.
Together with the atom at $x$, this proves~\eqref{eq:transition-kernel-Un}.
Since the right-hand side does not depend on~$n$, the kernel $Q_k$ is
time-homogeneous. It remains to verify the Markov property with respect
to the natural filtration $(\mathcal{G}_n)$.
Note that the conditional distribution
$\mathbb{P}(U_{n+1}\in\cdot\mid U_n=x)=Q_k(x,\cdot)$
was computed by conditioning only on $U_n=x$: the independence of
$X_{n+1}$ and $\mathcal{F}_n$ entered through $\mathbb{P}(X_{n+1}>y)=1-F(y)$,
while the conditional law of
$(X_{n-k+2:n},\dots,X_{n:n})$ given $U_n=x$ was supplied by
Lemma~\ref{lem:conditional-upper-order-stats}, which conditions
on $U_n$ alone and involves no other past information.
In particular, the regular conditional distribution of
$U_{n+1}$ given $U_n$ does not depend on $(U_k,\dots,U_{n-1})$, and
$(U_n)_{n\ge k}$ is a Markov chain with transition kernel~$Q_k$.
\end{proof}

\end{document}
